\chapter{Photometric (Camera) Systems}
\label{ch:camera}

\section{Imaging architectures}
\label{sec:camarch}

Photometric launch monitors capture a short burst of high-speed, IR-strobed
stereo images over a small capture volume beside (or above) the tee ---
roughly the first 30\,cm of ball flight --- and reconstruct ball and club
kinematics photogrammetrically. The commercial family tree:

\begin{itemize}
\item \textbf{Foresight Sports} defined the reference floor-unit design.
The GC2 (2010) used two cameras (``stereoscopic''), with club data via the
add-on HMT module (two more cameras). The GC3 / Bushnell Launch Pro (a
Foresight-built rebadge) use three cameras; the GCQuad and QuadMAX use
four cameras at the corners of the sensor window, capturing on the order
of 200 images of ball and club per shot at burst rates quoted up to
10,000\,fps; the ceiling-mounted GCHawk points the same quadrascopic
package downward~\cite{foresightgcquad,playbetterblp}. Onboard IR object
tracking detects the teed ball and the inbound clubhead and triggers the
burst; intrinsics/extrinsics are factory-calibrated (``no calibration
needed'')~\cite{foresightgcquad}.
\item \textbf{Uneekor} mounts overhead: EYE~XO/XO2 look straight down from
$\sim$3\,m with two (XO2: three) high-speed IR cameras at
3,000+\,fps~\cite{uneekoreyexo}. The overhead geometry sees the clubhead
and face region directly, enabling sticker-free club data and
impact-location video replay (Club Optix). The portable EYE~MINI reverts
to a side view and therefore requires club stickers.
\item \textbf{ProTee VX}: overhead, two synchronized high-speed cameras
with IR illumination and an ML shot-analysis pipeline; sticker-free club
and ball data including vertical/horizontal impact
point~\cite{proteevx}.
\item \textbf{Garmin Approach R50}: three high-speed cameras arranged
horizontally in a floor unit; purely optical ball and club measurement,
with clubhead data requiring reflective fiducials~\cite{garminr50}.
\item \textbf{SkyTrak / SkyTrak+}: the original SkyTrak was photometric and
ball-only; the SkyTrak+ added a dual Doppler radar module specifically for
club data while the cameras handle ball launch and spin~\cite{skytrakplus}.
\end{itemize}

\section{Ball measurement}

\subsection{Photogrammetric position and velocity}

With factory-known intrinsics $K_i$ and extrinsics $[R_i|\vect{t}_i]$,
each synchronized image set yields a 3D ball-center fix by triangulation:
back-project the matched detections as rays and solve the least-squares
intersection (or the DLT system $\tilde{\vect{x}}\times(P\vect{X})=0$),
refining by reprojection error. Ball speed is the displacement between
timestamped fixes; launch angle and direction are the components of the
initial velocity vector in the calibrated target frame. Two error sources
dominate: depth error grows quadratically with distance,
\begin{equation}
\label{eq:depth}
\sigma_Z = \frac{Z^2\,\sigma_{px}}{f\,B},
\end{equation}
for baseline $B$ and focal length $f$ --- which is why photometric units
keep the capture volume small and also exploit the known ball diameter
(21.34\,mm radius) as an auxiliary range cue; and center-finding bias ---
fitting the projected \emph{conic} outline rather than taking the blob
centroid avoids a systematic offset of up to $RZ/f$ pixels.

\subsection{Spin: dimple-pattern registration}
\label{sec:dimplespin}

The flagship photometric capability is markerless 3D spin. Successive
frames of the ball are registered on the sphere: the estimator searches
the rotation $R\in SO(3)$ that best maps the back-projected surface
texture (the dimple pattern, plus any logo or blemish) of frame $k$ onto
frame $k{+}1$ --- Foresight brands this \emph{spherical correlation};
Uneekor, \emph{Dimple Optix}~\cite{foresightspherical,uneekordimple}.
The recovered $R$ gives everything at once:
\begin{equation}
\label{eq:rotangle}
\hat{\vect{\omega}} = \mathrm{eig}_1(R), \qquad
\theta = \arccos\!\Bigl(\frac{\mathrm{tr}\,R - 1}{2}\Bigr), \qquad
\omega = \theta/\Delta t,
\end{equation}
i.e.\ the spin axis is the eigenvector of $R$ with unit eigenvalue and the
spin rate follows from the per-frame rotation angle. Reported accuracy is
1--3\% of true spin, independent of ball speed. The known failure mode is
rotational aliasing at high spin and low frame rate (the $n\cdot2\pi$
ambiguity), resolved by multi-hypothesis tracking or a
trajectory-consistency prior. Where resolution cannot support dimple
registration, systems fall back to \emph{marked balls}: Uneekor QED and
Rapsodo (RPT balls) track printed high-contrast features, giving closed-form
rotation from $\ge3$ correspondences via the Kabsch/Horn algorithm.

\section{Club measurement}
\label{sec:camclub}

\subsection{Fiducial markers: the Foresight approach}

Foresight floor units require retroreflective \textbf{fiducial dots} on
the clubface. The dot count maps directly to the data
tier~\cite{foresighthmt,foresightmarkers}: one dot suffices for club
speed, path, and attack angle (tracking a single 3D point through the
capture volume); \emph{four dots} --- placed on the vertical centerline,
equidistant from the horizontal centerline --- define the face plane and
enable face angle, dynamic loft, lie, closure rate
(\si{\degree\per\second}), and \textbf{impact location}. Mechanically: the
IR strobe makes the dots bloom in every frame; triangulating the dot
constellation across cameras and frames yields the 6-DOF pose of the face
through impact. Club speed is the pose translation rate; path and attack
angle are the velocity direction; face angle and dynamic loft are the
face-plane orientation at contact; impact location is the measured ball
position expressed in the face frame. The GC3/Launch Pro tier delivers
speed, path, attack, face angle, and impact via the same stickers but
omits dynamic loft, lie, and closure rate --- these remain
GCQuad/QuadMAX exclusives~\cite{playbetterblp}.

\subsection{Sticker-free overhead tracking}

Uneekor's overhead cameras segment the clubhead silhouette directly ---
no markers --- because the top-down view presents the head against the mat
with controlled IR illumination; the systems also record actual
high-speed impact video (face-on via Club Optix), from which impact
location is read~\cite{uneekoreyexo}. ProTee VX claims the same
sticker-free club set (speed, path, face, attack, dynamic loft, lie,
vertical and horizontal impact point) from its overhead ML
pipeline~\cite{proteevx}. The trade-off is installation: overhead units
need $\sim$3\,m ceilings, rigid mounting, and a floor-chart calibration.

\subsection{Calibration}

Factory-rigid multi-camera rigs ship calibrated; the user aligns only the
target line. Overhead systems require field calibration: Uneekor's
procedure places a printed calibration chart on the mat, levels it,
squares it to the screen, and aligns its crosshair to overlays on the
live camera feeds --- establishing the ground plane, hitting zone, and
target-line extrinsics~\cite{uneekorcalib}. This is the pattern a DIY
overhead build should copy.

\section{The indoor argument, and what cameras cannot do}

Photometric systems measure actual launch conditions --- including 3D spin
--- in the first foot of flight; nothing is lost when the ball hits a
screen 2.5\,m away. Radar indoors must infer from a truncated flight.
This is the structural reason camera units dominate indoor simulation
while radar dominates outdoor practice: each modality measures what the
other models~\cite{uneekorphotometric}. The camera's blind spots are
downrange (all flight is simulated from launch) and, for floor units, the
club face itself without stickers --- the same face-visibility problem
radar has, solved with markers instead of models.

\section{DIY and open-source photometric systems}
\label{sec:diy}

\textbf{PiTrac}~\cite{pitrac} is the flagship open-source photometric
build and the most relevant external project to OpenFlight. Instead of
multi-kfps cameras it uses \emph{strobed multi-exposure capture}: an IR
LED array pulses several times within one long exposure of a $\sim$\$50
Raspberry Pi Global Shutter camera, freezing multiple ball images in a
single frame. One camera watches the teed ball and detects launch; the
second captures the strobed flight images; a custom PCB sequences
trigger and strobe. Processing (C++/OpenCV): Hough-circle detection
locates ball images; strobe-timestamped positions give speed and launch
angles; 3D spin on all three axes is solved by searching candidate
rotations that best register the dimple imagery between exposures ---
Gabor-filter dimple enhancement is documented in the adjacent patent
literature~\cite{us12401909}. Total BOM $\approx$\$250--300. The
architecture is a direct descendant of the now-expired Wintriss patents
(\cref{sec:wintriss}) and of the GC2's flash-multi-exposure technique.

\begin{implication}
PiTrac proves that commodity global-shutter sensors + IR strobing deliver
photometric-grade launch and spin measurement at OpenFlight's price
point. For OpenFlight the natural division of labor is: keep the radar
core for speed/trigger robustness and outdoor use, and adopt a
PiTrac-style strobed camera as the spin/impact-location module indoors
--- the same fusion direction every commercial vendor has taken, from the
opposite starting modality.
\end{implication}
